Este capítulo apresenta os fundamentos teóricos necessários ao desenvolvimento do subsistema de cache, que são a metodologia de projeto de processadores, a linguagem Chisel, a arquitetura RISC-V, a hierarquia de memória e o funcionamento de memórias cache.

\section{Metodologia de projeto de processadores digitais}\label{sec:2.1}

A compreensão de um sistema computacional parte, na literatura de organização de computadores, de sua organização em níveis de abstração: cada nível oferece uma interface para o nível superior e oculta os detalhes do nível inferior~\cite{patterson2020computer}. Dessa maneira, essa hierarquia torna praticável o projeto de um processador, pois permite que o projetista raciocine sobre comportamento e temporização sem lidar, a cada decisão, com o comportamento físico dos transistores que compõem o circuito.

Sobre essa base de abstrações e de tecnologia se assenta o fluxo de projeto propriamente dito, que percorre estágios sucessivos até o circuito físico~\cite{weste2011cmos}. Nesse cenário, os primeiros passos são a especificação dos requisitos de área, desempenho e consumo e a modelagem funcional em linguagem de alto nível, com o intuito de validar preliminarmente o comportamento. Em seguida, passa-se à descrição estrutural em linguagem de descrição de hardware e verifica-se o circuito descrito por simulação em nível de transferência entre registradores, na qual bancadas de teste confrontam o comportamento obtido com o especificado~\cite{bergeron2006testbenches}. Uma vez verificado, o projeto é submetido à síntese lógica, que o traduz em uma lista de conexões de portas reais e à implementação física, em arranjo de portas programáveis em campo para prototipagem ou em circuito integrado de aplicação específica~\cite{weste2011cmos}.

Desse modo, o fluxo que orienta a organização do presente trabalho segue a mesma lógica, focando na fase de projeto lógico do ciclo de projeto de circuitos integrados, também chamada RTL\glsadd{RTL} (\emph{Register-Transfer Level}). Assim, o presente trabalho apresenta as formalidades necessárias de um sistema cache, a descrição feita em Chisel, a verificação funcional de cada bloco por ChiselTest e por simulação com Treadle e a síntese lógica com o auxílio da ferramenta Quartus Prime.

\section{Chisel como linguagem de descrição de hardware}\label{sec:2.2}

O Chisel (\emph{Constructing Hardware in a Scala Embedded Language}) não é uma linguagem de descrição de hardware no sentido convencional de Verilog e VHDL\glsadd{VHDL}, mas uma biblioteca de classes e funções em Scala que representa as primitivas necessárias para expressar circuitos digitais síncronos~\cite{schoeberl2025chisel}. Dessa distinção decorre a característica que organiza toda a relação entre código e circuito, em que um projeto em Chisel é um programa Scala que, ao ser executado, gera o circuito correspondente~\cite{schoeberl2025chisel}. Dessa maneira, não há correspondência direta entre linhas de código e portas lógicas, mas sim uma etapa de elaboração em que o programa constrói internamente a estrutura do hardware e constitui a descrição do circuito.

Em suma, o fluxo de ferramentas completa a passagem do código ao circuito. A classe Chisel é compilada pelo compilador Scala, junto com as bibliotecas Chisel e Scala, em um arquivo de classe executável pela máquina virtual Java. Sua execução por um driver gera uma representação intermediária flexível para RTL, sobre a qual o compilador FIRRTL\glsadd{FIRRTL} aplica transformações, até que o emissor de Verilog produza o código destinado à síntese~\cite{schoeberl2025chisel}. Nesse arranjo, o Verilog ocupa a posição de linguagem intermediária para teste e síntese.

Diante disso, o aspecto que mais diretamente sustenta a escolha do Chisel neste trabalho é a geração parametrizada de hardware. Por ser hospedada em uma linguagem de propósito geral com suporte a orientação a objetos e a programação funcional, a linguagem permite que características estruturais do circuito sejam expressas como parâmetros do gerador (argumentos de construtor, classes de configuração e funções com parâmetros de tipo) em vez de constantes repetidas no corpo da descrição~\cite{schoeberl2025chisel}. O efeito prático é que variações do mesmo circuito passam a ser instâncias distintas de um mesmo gerador, e não cópias modificadas do código. Para o objeto deste trabalho, isso significa que capacidade, tamanho de linha e grau de associatividade da cache podem ser tratados como argumentos, e que a avaliação experimental prevista no \refCap{4_Verificacao} se obtém por parametrização, sem reescrita da descrição.

\section{Arquitetura RISC-V}\label{sec:2.3}

A arquitetura do conjunto de instruções ocupa, conforme estabelecido na seção~\ref{sec:2.1}, a posição de contrato entre software e hardware. Desse contrato, o subsistema de memória herda um conjunto restrito de definições: o RV32I fixa a palavra em 32 bits, exige que instruções e dados residam em uma memória endereçável por byte e mantém a codificação das instruções em 32 bits fixos, independentemente do valor dos operandos~\cite{patterson2020computer}. Patterson e Hennessy apresentam essa última decisão como um compromisso explícito, formatos diferentes complicam a decodificação, mas permitem que todas as instruções tenham o mesmo tamanho~\cite{patterson2020computer}.

Todavia, o ponto em que a arquitetura determina diretamente o projeto da cache está nas instruções de acesso à memória. O RISC-V é uma arquitetura de carga e armazenamento, ou seja, apenas \textbf{load} e \textbf{store} acessam a memória, e as operações aritméticas nunca o fazem~\cite{patterson2020computer}. O conjunto básico oferece cinco variações de carga (\textbf{LB}, \textbf{LH} e \textbf{LW}, com extensão de sinal, e \textbf{LBU} e \textbf{LHU}, com extensão de zero) e três de armazenamento (\textbf{SB}, \textbf{SH} e \textbf{SW}), que gravam um byte, meia-palavra ou palavra~\cite{patterson2020computer}. O campo \textbf{funct3} da instrução seleciona simultaneamente o tamanho do acesso e o tipo de extensão~\cite{patterson2020computer}. Por consequência, impõe-se ao subsistema de memória duas obrigações: a escrita parcial de uma linha, que exige sinalização de quais bytes devem ser efetivamente gravados, e a leitura com seleção de tamanho e extensão de sinal, que precisa ocorrer entre o dado entregue pela cache e o registrador de destino.

Em suma, a cache retém um conjunto fechado de requisitos da arquitetura: palavra de 32 bits, espaço de endereçamento de 32 bits, memória endereçável por byte, acessos de byte, meia-palavra e palavra com extensão de sinal ou de zero, e alinhamento obrigatório. São esses aspectos que definirão os parâmetros da seção~\ref{sec:2.5.1}.

\section{Hierarquia de memória e princípio da localidade}\label{sec:2.4}

A memória de um sistema computacional não é um componente único, mas um conjunto de níveis organizados segundo um compromisso entre três grandezas que não podem ser maximizadas simultaneamente: capacidade, latência e custo por bit. Diante desse panorama, Patterson e Hennessy descrevem essa organização como uma pirâmide em que os níveis superiores são menores, mais rápidos e mais caros por bit armazenado, enquanto os inferiores são maiores, mais lentos e mais baratos~\cite{patterson2020computer}. No topo, situam-se os registradores do processador, logo abaixo, as memórias estáticas que compõem a cache, em seguida, a memória principal, construída com células dinâmicas, e, na base, o armazenamento secundário. A consequência prática dessa assimetria é que nenhuma tecnologia isolada oferece ao mesmo tempo a velocidade que o processador exige e a capacidade que os programas demandam.

Nesse sentido, a hierarquia responde a essa limitação por meio da replicação, ou seja, cada nível mantém cópias de parte do conteúdo do nível imediatamente inferior, de modo que uma requisição possa ser atendida pelo nível mais alto que disponha do dado~\cite{patterson2020computer}. Porém, se a eficácia da hierarquia dependesse apenas dessa replicação, o arranjo seria inútil, porque a memória principal é ordens de grandeza maior do que a cache e as cópias tenderiam a estar sempre ausentes. O que torna o arranjo viável é a existência de um padrão nos acessos, que constitui o princípio da localidade.

A localidade se manifesta em duas formas. A localidade temporal estabelece que um item acessado recentemente tem alta probabilidade de ser acessado novamente em um intervalo curto de tempo, o que decorre da estrutura dos programas tendo em vista que um laço executa repetidamente o mesmo conjunto de instruções~\cite{patterson2020computer}. Já a localidade espacial estabelece que itens cujo endereço é próximo ao de um item acessado recentemente também tendem a ser acessados, o que tem origem em dois fatos: as instruções são executadas sequencialmente, salvo em desvios e saltos, e as estruturas de dados, como vetores e matrizes, são armazenadas em posições adjacentes~\cite{patterson2020computer}. Assim, a existência desses dois padrões permite a uma memória pequena atender à maior parte das requisições de um programa.

\section{Memória cache}\label{sec:2.5}

A memória cache é o nível da hierarquia imediatamente abaixo dos registradores e tem por função manter cópias de parte do conteúdo da memória principal, de forma que a maior parte das requisições do processador seja atendida sem recorrer aos níveis inferiores. Sua eficácia não depende de armazenar os acessos mais frequentes, o que exigiria conhecer antecipadamente o comportamento do programa, mas de explorar a localidade descrita na seção~\ref{sec:2.4}. Para tal, a linha de cache, unidade de armazenamento e de transferência, agrega palavras adjacentes e permanece válida após o primeiro acesso, ao passo que tanto a repetição quanto a vizinhança dos acessos sejam aproveitadas~\cite{patterson2020computer}. Daí decorrem quatro decisões de projeto, conforme o enquadramento que Patterson e Hennessy aplicam à hierarquia em geral: onde um bloco pode ser colocado, como um bloco presente é identificado, qual bloco é substituído em caso de falta e o que ocorre em uma escrita~\cite{patterson2020computer}.

\subsection{Parâmetros arquiteturais}\label{sec:2.5.1}

A alternativa de uma memória única, que armazene instruções e dados sob o mesmo arranjo, herda da organização de von Neumann a economia de uma só estrutura e de um só conjunto de tags, e é a escolha natural em processadores que executam uma instrução por vez. Contudo, essa economia não se sustenta em um processador organizado em pipeline. Como a busca de instruções e o acesso a dados ocorrem em estágios distintos e simultâneos, uma memória única precisaria atender, no mesmo ciclo, uma leitura de instrução e uma leitura ou escrita de dado, o que exige duas portas de leitura e uma de escrita. Por isso, a separação em duas caches, uma para instruções e outra para dados, resolve o conflito pela raiz haja vista que cada uma atende uma única requisição por ciclo e, logo, precisa apenas de uma porta de leitura e de uma de escrita.

A especificação de uma cache se reduz à escolha de cinco parâmetros: a capacidade total, o tamanho da linha, o grau de associatividade e as políticas de escrita e de reposição~\cite{patterson2020computer}. Os dois últimos são tratados nas seções~\ref{sec:2.5.3} e~\ref{sec:2.5.4}, em razão do número de alternativas que cada um admite. Esta seção ocupa-se dos três primeiros, que definem a geometria da memória e, com eles, a decomposição do endereço apresentada na seção~\ref{sec:2.5.2}.

Nenhum dos três pode ser escolhido isoladamente, porque cada um desloca o compromisso entre três grandezas que se opõem: a taxa de acerto, o tempo de acerto e o custo em área. A título de exemplificação, aumentar a capacidade reduz as faltas de capacidade, mas amplia a área ocupada e tende a alongar o caminho combinacional de decodificação do endereço, o que eleva o tempo de acerto. Portanto, não existe configuração ótima em abstrato, sendo a escolha dependente do padrão de acesso esperado e do orçamento de área disponível no dispositivo de destino~\cite{patterson2020computer}.

O tamanho da linha incide diretamente sobre a localidade espacial. Nesse contexto, linhas maiores trazem mais palavras por transferência e tendem a reduzir a taxa de faltas em programas de acesso sequencial. Porém, há duas consequências contrárias. A primeira é que, para uma mesma capacidade, linhas maiores reduzem o número de linhas disponíveis, o que agrava as faltas de conflito. A segunda é que aumentam a penalidade de cada falta, porque mais palavras precisam ser transferidas antes que o dado solicitado seja entregue~\cite{patterson2020computer}.

A associatividade define quantas posições podem alojar um mesmo bloco da memória principal. No mapeamento direto, cada bloco tem uma única posição possível, determinada pelo resto da divisão do número do bloco pelo número de linhas da cache. Por outro lado, no mapeamento associativo por conjuntos, cada bloco pode ocupar qualquer uma das n vias de um conjunto determinado pelo índice~\cite{patterson2020computer}. Dessa maneira, o aumento da associatividade reduz as faltas de conflito, mas exige um comparador de tag por via, todos operando em paralelo, o que eleva simultaneamente o tempo de acerto e a área.

Sobre essa base incide uma restrição própria do dispositivo de destino. Os arranjos de portas programáveis em campo oferecem dois recursos distintos para armazenamento: os blocos de memória dedicados, cujas portas são sincronizadas e cuja leitura entrega o dado um ciclo após a apresentação do endereço, e a memória distribuída, construída com as tabelas de consulta da própria lógica programável, cuja leitura é combinacional~\cite{schoeberl2025chisel}. A escolha entre os dois não é livre, mas decorre do estilo de descrição empregado, e é ela que determina onde o arranjo cabe.

Contudo, os três parâmetros não são fixados na descrição, são argumentos do gerador. A geometria é expressa como uma classe de configuração, recebida pelo módulo da cache, pelo arranjo de dados, pelo controlador de reposição e pelo módulo de decomposição do endereço, de modo que as larguras de campo, o número de conjuntos, o número de vias e o número de palavras por linha derivam todos de um único ponto de definição. A consequência prática é que uma configuração distinta é uma instância distinta do mesmo gerador, e não uma cópia modificada do código, conforme a característica de parametrização discutida na seção~\ref{sec:2.2}.

\subsection{Organização e formato dos dados internos}\label{sec:2.5.2}

O endereço de 32 bits emitido pelo processador é decomposto em três campos, cujas larguras decorrem da geometria da configuração padrão, apresentada na seção~\ref{sec:3.2}. Nesse sentido, o deslocamento ocupa os 4 bits menos significativos, uma vez que a linha tem 16 bytes, e subdivide-se em dois subcampos: os 2 bits mais baixos selecionam o byte dentro da palavra, e os 2 seguintes selecionam a palavra dentro da linha. O índice ocupa os 7 bits seguintes, que endereçam os 128 conjuntos. A tag ocupa os 21 bits restantes. Dessa decomposição resulta uma propriedade que interessa ao desenho experimental: a largura da tag depende da capacidade e da associatividade, mas não do tamanho da linha. Como a capacidade é o produto entre o número de conjuntos, o número de vias e o tamanho da linha, a soma das larguras do índice e do deslocamento é constante para uma dada capacidade e associatividade, de modo que aumentar a linha reduz o índice na mesma medida em que amplia o deslocamento. Formalmente, a largura da tag é dada pela expressão (1).

tag = 32 $-$ log2(numSets) $-$ log2(lineBytes) (1)

Nessa perspectiva, cada linha reúne quatro campos, entregues pelo arranjo em um único agregado por via. O primeiro é o vetor de dados, composto por quatro palavras de 32 bits, o que totaliza 128 bits. O segundo é a tag, com os 21 bits acima. O terceiro é o bit de validade, que distingue uma linha com conteúdo legítimo de uma linha cujo conteúdo é resíduo de energização e não corresponde a endereço algum. O quarto é o bit de sujeira, que registra que a linha foi alterada e diverge da memória principal. Assim, a linha da cache de dados ocupa 151 bits, dos quais 128 são informação útil. Já a cache de instrução dispensa o bit de sujeira, por realizar exclusivamente leituras, e sua linha ocupa 150 bits.

Em suma, a geometria da configuração padrão (seção~\ref{sec:3.2}) se traduz em uma decomposição de 4, 7 e 21 bits, em uma linha de 151 bits por via e em um arranjo de 128 conjuntos. São esses números que definem o que a máquina de estados deve manipular e que determinam o custo em recursos da cache. Sobre esse arranjo incidem, agora, as decisões de política: o que ocorre em uma escrita e qual linha é descartada quando um conjunto está ocupado, o que constitui o objeto das seções~\ref{sec:2.5.3} e~\ref{sec:2.5.4}.

\subsection{Políticas de escrita e de alocação}\label{sec:2.5.3}

A política de escrita não é uma decisão única, mas duas independentes, que respondem a perguntas distintas. A primeira define o que ocorre quando uma operação de armazenamento encontra a linha já presente, isto é, em um acerto de escrita. A segunda define o que ocorre quando a operação não a encontra, em uma falta de escrita, e concerne a se a linha deve ser trazida para a cache antes que a gravação se efetive~\cite{patterson2020computer}. Como as duas perguntas admitem resposta em separado, as combinações resultantes são livres.

No acerto de escrita, a alternativa mais simples é a escrita direta, em que o dado é gravado simultaneamente na linha da cache e na memória principal, de modo que as duas cópias permanecem sempre idênticas. Por outro lado, a outra opção é a escrita de volta, em que o dado é gravado apenas na linha da cache, a qual é marcada como suja, e a memória principal só é atualizada quando essa linha for substituída~\cite{patterson2020computer}. A diferença entre as duas é de tráfego e de complexidade, mantendo-se a correção, haja vista que ambas preservam o valor observado pelo processador.

Na falta de escrita, a decisão é entre alocar e não alocar. Com alocação, a linha é trazida da memória principal e a gravação é aplicada sobre a cópia local, o que caracteriza a estratégia de leitura para posse. Por outro lado, sem alocação, a gravação é encaminhada diretamente à memória e a cache permanece inalterada, o que evita a transferência de uma linha cujo conteúdo será em parte descartado. A alocação costuma acompanhar a escrita de volta, e a não alocação costuma acompanhar a escrita direta, embora a combinação inversa também seja possível~\cite{patterson2020computer}. Então, o critério que orienta a escolha é o mesmo: alocar compensa quando se espera que a linha venha a ser reutilizada, e não compensa quando as escritas são isoladas e não se repetem.

\subsection{Política de reposição}\label{sec:2.5.4}

A política de reposição responde a uma pergunta que só se coloca quando o conjunto está cheio: qual das linhas presentes deve ceder lugar à que chega. Diferentemente da escrita, a qual define o que acontece com a cópia presente, se ela diverge da memória e em que momento essa divergência se resolve, a reposição define qual linha deixa a cache e, por consequência, se haverá despejo. Dessa distinção decorre que o custo de uma falta depende das duas decisões em conjunto, porque é a política de reposição que determina se a vítima está suja.

Em uma cache de mapeamento direto, essa pergunta não existe. Cada bloco tem uma única posição possível, de modo que a linha a substituir é sempre a mesma, e a política deixa de ser um parâmetro. A partir de duas vias, a escolha passa a influenciar a taxa de acerto, porque duas linhas legítimas podem competir pela mesma posição do conjunto, e é essa a razão pela qual a política figura entre os elementos de projeto ao lado da capacidade, do tamanho de linha e da associatividade~\cite{patterson2020computer}.

A primeira política disponível é a substituição da linha menos recentemente usada, chamada LRU\glsadd{LRU} (\emph{Least Recently Used}) que se apoia diretamente na localidade temporal descrita na seção~\ref{sec:2.4}: se um item acessado recentemente tende a ser acessado outra vez, então a linha cujo último acesso é o mais antigo é a menos provável de ser reutilizada, e é a que deve ceder lugar~\cite{patterson2020computer}. Para que essa política funcione, a cache precisa manter, por conjunto, um estado que registre a ordem de uso das vias. Esse estado é atualizado em todo acerto e é consultado apenas quando há falta com o conjunto cheio. A largura do estado cresce com a associatividade, e é por isso que implementações reais frequentemente empregam aproximações, como um campo de idade por via em lugar da ordem exata.

A segunda política é a de primeira a entrar, primeira a sair, denominada FIFO\glsadd{FIFO} (\emph{First In First Out}), que ignora os acertos e considera apenas a ordem de preenchimento, substituindo sempre a linha que está no conjunto há mais tempo~\cite{patterson2020computer}. Sua vantagem é o custo, porque dispensa o estado de recência e os comparadores que ele exige, resumindo-se a um contador por conjunto. Sua desvantagem é a insensibilidade ao uso tendo em vista que uma linha intensamente reutilizada pode ser descartada por ser antiga, ainda que seja a mais útil do conjunto.

A terceira política é a escolha aleatória da via, que serve como referência de comparação por não empregar nenhuma informação sobre o padrão de acesso. Essa escolha tem uma consequência favorável à verificação: como a sequência é reproduzível a partir do estado inicial do registrador, uma execução pode ser repetida exatamente, o que não ocorreria com uma fonte verdadeiramente aleatória. Assim, o mesmo teste produz a mesma sequência de vítimas, e uma divergência observada em uma execução pode ser investigada em outra.

\subsection{Consistência e FENCE.I}\label{sec:2.5.5}

Como cada cache, uma para instruções e outra para dados, mantém uma visão própria do conteúdo da memória e como nenhuma consulta a outra antes de responder, nada garante que uma instrução gravada pela cache de dados seja observada pela cache de instrução em uma execução subsequente. Patterson e Hennessy registram essa possibilidade como consequência direta da separação: a cache de instrução pode reter uma cópia que já não corresponde ao conteúdo da memória, e essa divergência só se resolve por intervenção explícita~\cite{patterson2020computer}. Os programas que gravam em posições que eles próprios executam dependem de que essa intervenção exista. Portanto, há uma ordem a garantir, que toda gravação anterior se torne visível a toda busca posterior.

Diante desse panorama, a FENCE.I é uma restrição de ordenação sobre o modelo de memória, que estabelece que as gravações realizadas antes dela sejam observáveis pelas buscas de instrução realizadas depois dela~\cite{waterman2019riscv}. A consequência dessa opção é que a obrigação recai integralmente sobre a implementação. Uma implementação em que a memória se comporte de forma suficientemente ordenada pode satisfazer a instrução sem executar nada; outra, como a deste trabalho, precisa traduzi-la em operações concretas sobre as caches. É justamente essa liberdade de mecanismo que a especificação preserva ao não o fixar.

Por fim, a solução é suficiente para o escopo deste trabalho e insuficiente para além dele. A FENCE.I ordena os dois fluxos dentro de um mesmo processador e nada estabelece sobre outros. Em um sistema com mais de um núcleo, cada um precisaria executar a instrução após receber uma interrupção do núcleo que gravou o código, porque o despejo e a invalidação locais não alcançam as caches dos demais~\cite{waterman2019riscv}. Como o processador em desenvolvimento é de núcleo único, essa extensão está fora do escopo delimitado na seção~\ref{sec:1.2}.
